\documentclass[11pt]{article}

\usepackage[margin=1in]{geometry}
\usepackage{amsmath,amssymb,amsthm}
\usepackage{enumitem}
\usepackage[colorlinks=true,linkcolor=blue,urlcolor=blue]{hyperref}

\newcommand{\TW}{Thatcher--Wright}
\newcommand{\Nat}{\mathbb N}
\newcommand{\powfin}{\mathcal P_{\mathrm{fin}}}
\newcommand{\Term}[1]{T_{#1}}
\newcommand{\Aut}{\mathcal A}
\newcommand{\bh}{\mathrm{bh}}
\newcommand{\MSO}{\ensuremath{\mathrm{MSO}}}
\newcommand{\enc}{\mathrm{enc}}

\title{Mathematical Issues and Review Notes for\\
  \texttt{Thatcher-Wright-expanding.tex}}
\author{}
\date{Review of the version at commit \texttt{fa1cd04}
  (\emph{expand\_all\_proof\_of\_tree\_automata}, July 4, 2026)\\
  Compiled: July 5, 2026}

\begin{document}
\maketitle

\section*{Scope and summary}

This report reviews the mathematical content of
\texttt{Thatcher-Wright-expanding.tex} against the source paper (Thatcher and
Wright, \emph{Generalized finite automata theory with an application to a
decision problem of second-order logic}, Math.\ Systems Theory 2 (1968),
57--81; \texttt{Thatcher-Wright.pdf} in this directory).  Line numbers refer to
\texttt{Thatcher-Wright-expanding.tex} at the commit named above.

The overall verdict is positive.  The expanded proofs of Theorems~1--4,
Lemmas~5--7, Lemma~10, Lemma~11, Lemma~12, Lemma~13, Lemmas~15--16, Lemma~18,
Lemma~19, Theorem~17, the first-order elimination, the Doner decidability
proof, and the derivation of lecture-note Theorem~26 were checked in detail and
are correct.  Two findings from the previous review round (June 2026 review of
the pre-expansion version) have been repaired, in both cases by deliberately
\emph{correcting} the source paper; these deviations are documented in
Section~\ref{sec:resolved} and should be kept.

The findings of the present round are:
\begin{enumerate}[label=\textbf{C\arabic*.},leftmargin=*]
  \item the induction step in the expanded proof of the analysis theorem
    (TW Theorem~9) is invalid as written: it restores a single removed subterm
    using the complex product \(\cdot_\lambda\), which by definition
    substitutes at \emph{every} \(\lambda\)-occurrence.  The statement being
    proved is true; the proof step is not.  A repair is given in
    Section~\ref{sec:critical}.
\end{enumerate}
The remaining findings (Sections~\ref{sec:obligations}--\ref{sec:path}) are not
errors but load-bearing lemmas that the text uses without proof, plus one
design decision that the added derivation of Theorem~26 forces on the Lean
formalization.

\section{Resolved findings and deliberate corrections of \TW{}}
\label{sec:resolved}

The June 2026 review of the pre-expansion draft reported two mathematical
errors inherited from the source paper.  Both are now repaired.  Because the
repairs silently deviate from \TW{}, we record them here so that later readers
comparing against the paper do not mistake the corrections for transcription
errors.  We recommend adding a short remark to the main file at each location.

\subsection{TW Lemma 19(1): the literal conjunction equation is false}

\TW{} (p.~79--80) prove closure of recognizable relations under conjunction by
lifting a relation \(R(\alpha,\beta)\) to
\(R'(\alpha,\beta,\gamma)\Leftrightarrow R(\alpha,\beta)\) via
\[
  \widehat\pi^{-1}(e(R))\cap U_3=e(R'),
\]
where \(\pi\) deletes the third track.  This equation is \emph{false} for
canonical encodings.  Counterexample: take \(R=\{(\emptyset,\emptyset)\}\), so
\(e(R)=\{\lambda\}\), and consider
\[
  t=e(\emptyset,\emptyset,\{1\})
   =(0,0,0)\bigl((0,0,1)(\lambda,\lambda),\;\lambda\bigr).
\]
Then \(t\in e(R')\cap U_3\), but
\(\pi(t)=(0,0)((0,0)(\lambda,\lambda),\lambda)\) contains the forbidden
subterm \((0,0)(\lambda,\lambda)\), is therefore not canonical, and so does not
lie in \(e(R)\).  Hence \(t\notin\widehat\pi^{-1}(e(R))\).  In general, a
deleted-track projection of a canonical term need not be canonical, because the
remaining tracks may not need the whole tree shape.

The current file repairs this correctly (lines 2267--2338): it first
\emph{saturates} \(e(R)\) by grafting all-zero trees,
\(e(R)\cdot_\lambda Z_k\) with \(Z_k=\{\mathbf 0_k(\lambda,\lambda)\}^{*_\lambda}\),
and only then pulls back and intersects with \(U_m\):
\[
  L_I=\pi_I^{-1}\bigl(e(R)\cdot_\lambda Z_k\bigr)\cap U_m .
\]
The counterexample above was re-checked against this definition and passes
(\(\pi(t)\) is an all-zero term, hence lies in
\(\{\lambda\}\cdot_\lambda Z_2\)).  The lifting fact and its use in the
conjunction case are correct.  \textbf{Keep this construction; do not
``simplify'' it back to the paper's version.}

\subsection{TW Lemma 12: the paper's construction misses two edge cases}

\TW{}'s automaton for \(U^{*_\delta}\) (p.~71) keeps the state set \(A\), sets
\(\rho_c=\{\alpha_c\}\) for all constants, and uses \(A_F\) as final set.  This
fails on two inputs:
\begin{itemize}[leftmargin=*]
  \item the bare term \(\delta\) lies in \(U^{*_\delta}\) by definition
    (\(X_0=\{\delta\}\)), but the paper's automaton accepts it only when
    \(\delta\in U\);
  \item a constant \(c\in U\) can never collapse to the \(\delta\)-role,
    because the collapse clause exists only at positive-arity transitions, so
    terms of the form \(u[\delta:=c]\) are wrongly rejected.
\end{itemize}
The current file's construction (lines 1351--1553) repairs both: it uses a
tagged copy \(\widehat A\) plus a fresh state \(q_\lambda\), makes
\(q_\lambda\) initial at the constant \(\lambda\), adds the initial clause
``\(q_\lambda\) if \(c\ne\lambda\) and \(\alpha_c\in A_F\)'', and accepts with
final set \(\{q_\lambda\}\).  The four invariants (2), (3), (Q), (A) were
checked and the proof is correct.  This is again a deliberate correction of the
source paper and should be kept.

\subsection{Other items now in order}

The \(\cup\{\lambda\}\) case of Lemma~15 (encoding of the empty set, line
1952), the zero-track sentence case \(U_0=\{\lambda\}\) in Lemma~19 and in the
Doner proof, the explicit disjointness assumption in Lemma~11, and the
rewritten minimal-size proof of Lemma~6 (lines 839--875, which replaces \TW{}'s
iterative shrinking by a well-founded minimal-witness argument and is the
preferable form for Lean) are all resolved or improved relative to the previous
round.

\section{Critical finding: the analysis-theorem induction step}
\label{sec:critical}

\paragraph{Location.} TW Theorem~9 (analysis theorem), converse inclusion of
claim \((*)\), lines 1056--1089; the invalid step is lines 1081--1087.

\paragraph{Statement being proved.} With
\(X_0=\{\lambda_{k+1}\}\), \(X_{r+1}=X_r\cup(V_{k+1}\cdot_{\lambda_{k+1}}X_r)\),
the text proves by induction on the number \(r\) of non-leaf transitions to
\(a_{k+1}\) that
\[
  t\in T_j^{k+1}
  \;\Longrightarrow\;
  t\in T_j^k\cdot_{\lambda_{k+1}}X_r .   \tag{\(**\)}
\]

\paragraph{The problem.} In the inductive step, a lowest subterm occurrence
\(v\in V_{k+1}\) is replaced by \(\lambda_{k+1}\), giving \(t_0\), and the
induction hypothesis places \(t_0\in T_j^k\cdot_{\lambda_{k+1}}X_r\).  The text
then concludes
\[
  t\in\bigl(T_j^k\cdot_{\lambda_{k+1}}X_r\bigr)
        \cdot_{\lambda_{k+1}}V_{k+1}.
\]
This does not follow.  The complex product \(\cdot_{\lambda_{k+1}}\)
substitutes at \emph{every} occurrence of \(\lambda_{k+1}\), whereas the
restoration must substitute at exactly \emph{one} occurrence (the position of
\(v\)) and leave all other \(\lambda_{k+1}\)-leaves of \(t_0\) unchanged.
Since every element of \(V_{k+1}\) has a positive-arity root pattern from
\(U_{k+1}\), we have \(\lambda_{k+1}\notin V_{k+1}\), so the other leaves
cannot be preserved by choosing \(\lambda_{k+1}\) itself from \(V_{k+1}\).

\paragraph{Counterexample to the step.} Consider any automaton containing a
binary symbol \(g\) with \(\alpha_g(a_{k+1},a_{k+1})=a_j\), \(j\le k\), and a
pattern \(v=f(\lambda_{i_1},\ldots,\lambda_{i_r})\in U_{k+1}\) none of whose
leaves is \(\lambda_{k+1}\) (recall \(U_{k+1}\subseteq V_{k+1}\), since
\(\lambda_i\in T_i^k\) for every \(i\)).  Let
\[
  t=g\bigl(v,\lambda_{k+1}\bigr).
\]
Then \(t\in T_j^{k+1}\) with exactly one non-leaf transition to \(a_{k+1}\),
and \(t_0=g(\lambda_{k+1},\lambda_{k+1})\in T_j^k\).  But every element of
\(T_j^k\cdot_{\lambda_{k+1}}V_{k+1}\) has \emph{both} holes of \(t_0\) filled
with \(V_{k+1}\)-terms, and no element of any \(T_j^k\)-term substitution can
produce the surviving right-hand leaf \(\lambda_{k+1}\).  Hence
\(t\notin(T_j^k\cdot_{\lambda_{k+1}}X_0)\cdot_{\lambda_{k+1}}V_{k+1}\),
although \((**)\) itself is true for \(t\):
\(t=g(\lambda_{k+1},\lambda_{k+1})\cdot_{\lambda_{k+1}}\{v,\lambda_{k+1}\}\)
using the per-hole independence of the product, so
\(t\in T_j^k\cdot_{\lambda_{k+1}}X_1\).

For historical accuracy: \TW{}'s own proof of their statement (III) contains
the same one-occurrence-versus-all-occurrences slip, so the expansion is a
faithful transcription of a flawed argument.  The claim \((*)\),
\(T_j^{k+1}=T_j^k\cdot_{\lambda_{k+1}}W_{k+1}\), is true.

\paragraph{Secondary defects of the same passage.}
\begin{itemize}[leftmargin=*]
  \item Even granting the restoration, the final inclusion
    \((T_j^k\cdot X_r)\cdot V_{k+1}\subseteq T_j^k\cdot X_{r+1}\) needs
    \(X_r\cdot_{\lambda}V\subseteq X_{r+1}\), while the stages are defined with
    \(V\) on the \emph{left}, \(X_{r+1}=X_r\cup(V\cdot_\lambda X_r)\).  The
    needed inclusion is provable, but only via same-symbol associativity,
    monotonicity, and its own induction; it is not immediate.
  \item The induction variable ``number of non-leaf transitions to
    \(a_{k+1}\) occurring in the computation of \(t\)'' has no formal
    definition in the file.  In Lean it would require instrumenting the run
    with the multiset of produced states --- avoidable machinery.
\end{itemize}

\paragraph{Recommended repair (removes all three defects at once).}
Replace the transition-counting induction by a \emph{topmost decomposition}
and structural induction:
\begin{enumerate}[leftmargin=*]
  \item Given \(t\in T_j^{k+1}\), mark the \emph{topmost} subterm occurrences
    of \(t\) whose root state is \(a_{k+1}\) (these are pairwise incomparable).
    Replacing all of them \emph{simultaneously} by \(\lambda_{k+1}\) yields a
    skeleton \(u\).  By maximality, no internal node of \(u\) produces
    \(a_{k+1}\), so \(u\in T_j^k\); moreover every
    \(\lambda_{k+1}\)-leaf of \(u\) is a marked position (an original
    \(\lambda_{k+1}\)-leaf of \(t\) has state \(a_{k+1}\) and hence lies in a
    topmost marked occurrence).
  \item Show that each marked subterm \(w\) lies in
    \(W_{k+1}=V_{k+1}^{*_{\lambda_{k+1}}}\), by strong induction on the size of
    \(w\): either \(w=\lambda_{k+1}\in W_{k+1}\), or the root of \(w\) is a
    non-leaf transition to \(a_{k+1}\), in which case \(w\) is a
    \(U_{k+1}\)-pattern whose holes carry terms handled by the same
    decomposition, giving \(w\in V_{k+1}\cdot_{\lambda_{k+1}}W_{k+1}\subseteq
    W_{k+1}\).  The last inclusion is the max-stage argument already used at
    (1) of Lemma~12 (lines 1405--1420) and in Lemma~15.
  \item Conclude \(t\in T_j^k\cdot_{\lambda_{k+1}}W_{k+1}\) in a single
    application of the product, with the marked subterms as the per-hole
    choices.  No occurrence counting and no one-at-a-time restoration is
    needed.
\end{enumerate}
This is also the form we recommend for the Lean proof, should the analysis
direction ever be formalized (note that the Courcelle pipeline does not need
it; see Section~\ref{sec:path}).

\section{Load-bearing lemmas used without proof}
\label{sec:obligations}

None of the following is an error, but each is an unproven statement that
multiple expanded proofs rely on.  In a formalization each becomes a named
obligation; we list them with their use sites and suggested treatment.

\subsection{The canonical-form characterization}
\label{sec:canonical}

\paragraph{Location.} Definitions at lines 1912--1917 and 2025--2030
(``the smallest term in \(c^{-1}(A)\), equivalently the \emph{unique}
representative with no subterm \(\mathbf 0(\lambda,\lambda)\)'').

\paragraph{Used by.} Both inclusions of Lemmas~15 and~16, the subset and
successor arguments of Lemma~18, claim (2) of the lifting fact in Lemma~19,
the recovery of \(\bar B\) from \(q\in U_{m-1}\) in the existential case, and
the identity \(\iota(\enc(T))=e(D_T,(P_a^T))\) in the derivation of
Theorem~26 (lines 122--128).

\paragraph{What must be proved.} Existence and uniqueness: every tuple
\(\bar A\) has exactly one \(c_n\)-preimage without the forbidden subterm, and
the order-theoretic ``smallest'' description agrees with it.  Recommended
treatment for Lean: do not define \(e\) through an order at all.  Define
\(e(\bar A)\) by recursion on the finite tuple (root label from
\((\epsilon\in A_i)_i\); recurse on the shifted tuples; return \(\lambda\) when
all components are empty), then prove
\begin{enumerate}[label=(\roman*),leftmargin=*]
  \item \(c_n(e(\bar A))=\bar A\);
  \item \(e(\bar A)\) has no forbidden subterm;
  \item any term without a forbidden subterm equals \(e\) of its decode.
\end{enumerate}
The key auxiliary fact for (iii) is:
\emph{a term whose decode is the all-empty tuple and which has no forbidden
subterm is \(\lambda\)} (an all-zero-labelled term of depth \(\ge2\) must
contain \(\mathbf 0(\lambda,\lambda)\) at the bottom).  This one lemma settles
the \(\lambda\)-versus-non-\(\lambda\) boundary in the uniqueness induction.

\subsection{The reduce/graft lemma (claim (2) of the lifting fact)}

\paragraph{Location.} Lines 2317--2332; re-used at lines 2402--2409.

\paragraph{Statement.} For canonical \(t=e(A_1,\ldots,A_m)\) and a coordinate
selection \(I\),
\[
  \pi_I(t)\in\{e(A_{i_1},\ldots,A_{i_k})\}\cdot_\lambda Z_k .
\]
The file proves this by a two-line sketch.  It is the real content of the
lifting fact: an arbitrary (possibly non-canonical) term is its own canonical
reduct with all-zero trees grafted at \(\lambda\)-leaves.  It deserves to be a
named, free-standing lemma (``every term of \(\Term{\Sigma_k}\) lies in
\(\{e(c_k(t))\}\cdot_\lambda Z_k\)''), proved by structural induction using the
zero-decode fact of Section~\ref{sec:canonical}.  Two further facts are used
silently nearby: grafting \(Z_k\)-terms does not change \(c_k\), and
\(Z_k=\{\mathbf 0_k(\lambda,\lambda)\}^{*_\lambda}\) equals the set of all
all-zero terms (asserted at lines 2279--2283 with a one-line proof note).

\subsection{An algebra of complex products}

The expanded proofs use, without statement or proof:
\begin{itemize}[leftmargin=*]
  \item \emph{same-symbol associativity}
    \((U\cdot_\lambda V)\cdot_\lambda T=U\cdot_\lambda(V\cdot_\lambda T)\):
    Lemma~18's converse (composing chain wrappers, lines 2246--2253), the
    secondary inclusion in Theorem~9, and Lemma~19;
  \item \emph{mixed-symbol rearrangement} for the nested products
    \(U_{k+1}\cdot_{\lambda_1}T_1^k\cdots\cdot_{\lambda_m}T_m^k\)
    (lines 1026--1033) and
    \(L_F\cdot_{\lambda_1}C_1\cdots\cdot_{\lambda_m}C_m\)
    (lines 1107--1121).  Caution: \TW{} themselves note that mixed-symbol
    associativity \emph{fails} in general
    (\(\{t\}=(\{\delta\}\cdot_\lambda\{\lambda\})\cdot_\delta\{t\}
      \ne\{\delta\}\cdot_\lambda(\{\lambda\}\cdot_\delta\{t\})=\{\delta\}\)).
    The uses here are sound because the inserted sets contain no occurrence of
    the other hole symbols, but that side condition must be made explicit;
  \item \emph{continuity / max-stage}: \(U\cdot_\lambda\bigl(\bigcup_rX_r\bigr)
    =\bigcup_r(U\cdot_\lambda X_r)\) for increasing chains --- stated and used
    correctly three times (Lemma~12 (1), Lemma~15, Lemma~16), but worth
    isolating once;
  \item the composition inclusion \(X_r\cdot_\lambda V\subseteq X_{r+1}\)
    discussed in Section~\ref{sec:critical}.
\end{itemize}
Recommendation: develop these once as a small substitution algebra
(associativity for a fixed hole symbol, conditional commutation for distinct
hole symbols, monotonicity, continuity, the \(\cdot_\lambda\emptyset\)
filter), and let every proof cite it.  Ad hoc re-derivation inside each proof
is exactly the pattern that produces oversized Lean proofs.

\subsection{Species extension and reduct transfer}

\paragraph{Location.} Theorem~9 replaces the constants of \(\Sigma\) by fresh
ones (\(\Sigma^+\), lines 965--978); Lemmas~15, 16, 18 add a temporary
\(\delta\); Theorem~8's synthesis direction says recognizability is preserved
``even after adding finitely many fresh constants'' (lines 1129--1140).

\paragraph{Missing statement.} If \(U\subseteq\Term{\Sigma}\) and \(U\) is
recognizable over an extension \(\Sigma'\supseteq\Sigma\) by fresh constants,
then \(U\) is recognizable over \(\Sigma\) (restrict the automaton to the
reduct; runs of \(\Sigma\)-terms never consult the new constants).  Trivial on
paper, but in a formalization terms over different species live in different
types, so one needs signature-inclusion homomorphisms
\(\Term{\Sigma}\hookrightarrow\Term{\Sigma'}\) and transfer lemmas for \(h\)
and \(\bh\).  Design options: (i) a graded-signature functorial interface in
which inclusions are ordinary signature maps and Theorem~4's machinery is
reused, or (ii) one fixed ambient species with the set of admissible constants
as a parameter, avoiding type-level extension entirely.

\subsection{Effectiveness of truncation}

\paragraph{Location.} Lemma~13 defines \(A_V=\{h_\Aut(v):v\in V\}\)
set-theoretically (lines 1578--1583); the effectiveness paragraphs of
Theorem~17 (lines 2490--2493) and the Doner proof claim every step is
effective.

\paragraph{Gap.} For arbitrary \(V\) the set \(A_V\) is not computable, so
Lemma~13 as stated is a pure existence result.  Effectiveness holds in every
actual use because the only \(V\) employed is \(Z_{m-1}\), which is given by an
explicit automaton; for automaton-given \(V\), \(A_V\) is computable by a
reachable-state analysis (the same induction as Theorem~7).  One remark in the
main file would close this; a computable Lean \texttt{compile} function must
implement it.

\subsection{Standing assumptions and small elisions}

\begin{itemize}[leftmargin=*]
  \item The species definition (lines 170--181) still does not carry \TW{}'s
    standing assumptions ``\(\Sigma\) finite'' and
    ``\(\Sigma_0\ne\emptyset\)''.  Finiteness is invoked inside Theorem~7
    (line 899); it should be a blanket hypothesis of the section.
  \item Lemmas~15/16, right-to-left inclusion: the argument that substitution
    creates no forbidden subterm at the junction silently uses
    \(\lambda\notin W_n\) (every closure element is \(\delta\) or has a
    binary root), so a pattern child slot filled by an inserted term is never
    \(\lambda\).  True, but state it; in Lean it is an explicit lemma.
  \item Derivation of Theorem~26: the relativization equivalence (1)
    (lines 140--145) is asserted ``by induction on \(\Phi\)'' --- this
    induction over full \MSO{} syntax is a genuine proof obligation, as are
    the weak-\MSO{} definability of the well-formedness conditions
    (prefix-closedness of \(D\), partition of \(D\) by the \(P_a\)) and the
    empty-tree edge case \(\enc(T)=\bot\), \(\iota(\bot)=\lambda\).
\end{itemize}

\section{Design decision forced by the new Theorem-26 derivation}
\label{sec:path}

The added derivation of lecture-note Theorem~26 (lines 94--166) is
mathematically correct, but it commits to the \emph{literal \(N_2\) route}:
its dependency cone is
\[
  \text{Thms 1--4}
  \;+\;\text{regularity (Lemmas 10--12, Thm 8)}
  \;+\;\text{truncation (Lemma 13)}
  \;+\;\text{encodings (Lemmas 15--16)}
  \;+\;\text{Lemmas 18--19}
  \;+\;\text{FO elimination}
  \;+\;\text{relativization},
\]
together with all obligations of Section~\ref{sec:obligations}.  Meanwhile the
closing section of the file (lines 2652--2656) still recommends the
\emph{track route} for the Courcelle formalization: variable assignments as
extra Boolean tracks on the fixed input tree, existential quantification as
plain projection.  On the track route the tree shape is shared by all tracks,
so no canonicalization exists at all: negation is plain complementation,
projection introduces no redundant subtrees, and the entire canonical/reduce/
graft/truncation apparatus (Sections~\ref{sec:canonical}--4.3, Lemmas~12--13,
15--16, and the \(Z_k\) machinery of Lemma~19) is unnecessary.  Its dependency
cone is only: Theorems~1--4, atomic track automata, the formula induction, and
a relativization-free semantics transfer.

Both routes cannot be ``the'' formalization target simultaneously.  Our
recommendation, unchanged from the previous review: formalize the track route
for Courcelle, and keep the \(N_2\)-literal derivation (and this file's
Sections~5--10) as paper-level provenance.  If the \(N_2\) route is chosen
instead, the obligations of Section~\ref{sec:obligations} move onto the
critical path, and the Theorem~9 repair of Section~\ref{sec:critical} is still
\emph{not} needed (only the synthesis direction of Theorem~8 is used).

\section{Notes on the linear-time accounting}

The planned treatment of ``linear time'' --- an instrumented evaluator that
increments a counter per step, with the final theorem bounding the counter by
a linear function of the input size --- fits the revised text well:
\begin{itemize}[leftmargin=*]
  \item the padded-term size bound \(|\enc(T)|\le2|T|+1\) is now explicit
    (line 165) and is exactly the encoding-size lemma the cost argument needs;
  \item the run is described as ``one transition per symbol of
    \(\enc(T)\)'' (lines 163--166), matching a one-tick-per-node cost
    convention;
  \item the derivation cleanly separates the compile phase (inverse
    projection, determinization --- cost uncontrolled, absorbed into the
    \(f(|\Phi|)\) constant) from the run phase.  The Lean statement should
    preserve this separation: prove the cost bound only for the run of the
    final deterministic automaton, with a constant depending on the automaton
    but not on the input tree, alongside an agreement lemma tying the
    instrumented evaluator to \(h_\Aut\).
\end{itemize}

\section{Verified coverage}

For the record, the following were checked in detail against the source paper
and found correct in the reviewed version: the derivation of Theorem~26 from
\TW{} (modulo the obligations listed above); the full proofs of Theorems~1,
2, 3, 4 and Lemmas~5, 6 (improved), 7; Lemma~10's singleton automaton;
Lemma~11 including both invariants; Lemma~12's corrected \(q_\lambda\)
construction and all four invariants; Lemma~13; the set-only translation and
its verification; Lemmas~15 and~16 including the \(\cup\{\lambda\}\) case;
Lemma~18 (both the \((1,0)\)-avoidance automaton and the chain construction
\(S_i=H^{*_\delta}\cdot_\delta\{B_i\}\), whose base terms and wrappers were
checked to encode and preserve the successor relation and canonicity);
Lemma~19's \(Z_k\)-saturated lifting fact, negation, and existential case
including the \(m-1=0\) degenerate alphabet; the assembly of Theorem~17; and
the Doner decidability proof including the zero-track sentence test.  The sole
exception is the analysis theorem's converse induction step
(Section~\ref{sec:critical}).

\end{document}
